%% main.tex — Empirical Software Engineering (Springer Nature)
%% Repository: behavioral-fsm-benchmark
%% Results populated only after campaign freeze (M9)

\documentclass[pdflatex,sn-mathphys-num]{sn-jnl}

\usepackage[T1]{fontenc}
\usepackage[utf8]{inputenc}
\usepackage{microtype}
\usepackage{graphicx}
\usepackage{booktabs}
\usepackage{multirow}
\usepackage{array}
\usepackage{float}
\usepackage{placeins}
\usepackage{enumitem}
\usepackage{amsmath,amssymb}
\usepackage{url}
\usepackage{xurl}
\usepackage[hypertexnames=false]{hyperref}

\raggedbottom

\title[Beyond Structural Validity]{%
  Beyond Structural Validity: Evaluating Behavioral Correctness, Robustness, and
  Reproducibility of LLM-Generated Finite State Machines from Natural-Language
  Requirements%
}

\author*[1]{\fnm{C\'esar Andr\'es} \sur{S\'anchez}}
\affil*[1]{%
  \orgdiv{CRIA-BDHS Research Group},%
  \orgname{Higher Polytechnic School of Technology and Science, Camilo Jos\'e Cela University},%
  \orgaddress{%
    \street{Calle Castillo de Alarc\'on 49},%
    \city{Villafranca del Castillo, Madrid},%
    \postcode{28692},%
    \country{Spain}%
  }%
}

\abstract{%
\textbf{Context.}
Practitioners increasingly use large language models (LLMs) to draft finite state machines (FSMs)
from natural-language requirements for model-based testing and simulation.
Prior work on local open-weight models established that many outputs are structurally valid
yet often fail strict determinism checks; behavioral correctness, robustness under requirement
change, and multi-run reproducibility remain insufficiently measured.
\textbf{Objective.}
We empirically investigate whether structurally valid LLM-generated FSMs behave as specified
when evaluated with approved gold references, executable behavioral test suites, and guard-aware
determinism analysis.
\textbf{Method.}
We extend the FSM-Bench-20 benchmark (20 systems, local Ollama protocol) with tiered gold FSMs
(12 core, 8 stretch, 3 pilot), a decidable guard DSL, and an offline evaluation framework
reporting structural gates (G1--G3a), behavioral test-suite agreement, gold conformance,
coverage, perturbation deltas, and reproducibility variance.
Four frozen campaigns address seven pre-registered research questions on the structural--behavioral
gap, failure taxonomy, proxy validity, robustness, and metric stability.
\textbf{Results.}
Reported after campaign freeze.
\textbf{Conclusions.}
Reported after campaign freeze.
}

\keywords{finite state machines, large language models, behavioral correctness, model-based testing, robustness, reproducibility, empirical software engineering}

\begin{document}

\maketitle

\section{Introduction}
\label{sec:introduction}

%% Goal: Motivate the structural--behavioral gap; position vs FSM-Bench-20 (IST 2026);
%% state contributions C1--C6; outline paper structure (§\ref{sec:study-design}--§\ref{sec:conclusion}).

\subsection{Motivation and problem}
\label{sec:introduction:motivation}

\subsection{Research gap}
\label{sec:introduction:gap}

\subsection{Contributions}
\label{sec:introduction:contributions}

\subsection{Paper organization}
\label{sec:introduction:organization}

\section{Related Work}
\label{sec:related-work}
%% Objective: LLM specification synthesis, MBT benchmarks, robustness, reproducibility; positioning vs FSM-Bench-20.

\section{Study Design}
\label{sec:study-design}

%% Goal: Present RQ1--RQ7, hypotheses H1--H6, experimental factors, campaigns C0--C4,
%% and analysis plan. Authority: docs/study_design.md, docs/evaluation_protocol.md.

\subsection{Research questions}
\label{sec:study-design:rq}

\subsection{Hypotheses}
\label{sec:study-design:hypotheses}

\subsection{Independent and dependent variables}
\label{sec:study-design:variables}

\subsection{Experimental factors and campaigns}
\label{sec:study-design:campaigns}

\subsection{Analysis plan}
\label{sec:study-design:analysis}

\subsection{Validity overview}
\label{sec:study-design:validity}

\section{Benchmark: \texorpdfstring{\textsc{FSM-Bench-Next}}{FSM-Bench-Next}}
\label{sec:benchmark}
%% Objective: Reference FSMs, test suites, guard semantics, metrics, pipeline. Authority: paper/notes/benchmark_specification.md

\subsection{Relationship to \texorpdfstring{\textsc{FSM-Bench-20}}{FSM-Bench-20}}
\label{sec:benchmark:upstream}

\subsection{System tiers}
\label{sec:benchmark:tiers}

\subsection{Reference FSMs and test suites}
\label{sec:benchmark:artifacts}

\subsection{Evaluation layers and metrics}
\label{sec:benchmark:metrics}

\subsection{Scoring}
\label{sec:benchmark:scoring}

\section{Experimental Setup}
\label{sec:experimental-setup}
%% Objective: Models, campaigns C0--C4, baseline import, provenance.

\subsection{Models and inference}
\label{sec:experimental-setup:models}

\subsection{Structural baseline}
\label{sec:experimental-setup:baseline}

\subsection{Campaigns}
\label{sec:experimental-setup:campaigns}

\subsection{Provenance and replication}
\label{sec:experimental-setup:provenance}

\section{Results}
\label{sec:results}
%% Objective: One subsection per RQ (§6.1--§6.7). Populate after campaign freeze only.

\subsection{Structural--behavioral gap (RQ1)}
\label{sec:results:gap}

\subsection{Behavioral failure taxonomy (RQ2)}
\label{sec:results:failures}

\subsection{Proxy metrics and gold conformance (RQ3)}
\label{sec:results:proxies}

\subsection{Overall perturbation sensitivity (RQ4)}
\label{sec:results:perturbation-overall}

\subsection{Perturbation-type effects (RQ5)}
\label{sec:results:perturbation-types}

\subsection{Run-to-run variance (RQ6)}
\label{sec:results:variance}

\subsection{Metric stability (RQ7)}
\label{sec:results:stability}

\subsection{Integrated summary}
\label{sec:results:summary}

\section{Discussion}
\label{sec:discussion}

%% Goal: Interpret RQ1--RQ7 answers; implications for MBT, RE practice, and benchmark design;
%% relation to IST descriptive findings (background only, no new EMSE numbers pre-freeze).

\subsection{Structural validity versus behavioral correctness}
\label{sec:discussion:gap}

\subsection{Failure modes and test-suite design}
\label{sec:discussion:failures}

\subsection{Construct validity of proxy metrics}
\label{sec:discussion:proxies}

\subsection{Robustness and requirements evolution}
\label{sec:discussion:robustness}

\subsection{Reproducibility and benchmark practice}
\label{sec:discussion:reproducibility}

\subsection{Practical recommendations}
\label{sec:discussion:recommendations}

\subsection{Limitations}
\label{sec:discussion:limitations}

\section{Threats to Validity}
\label{sec:threats}

%% Goal: Internal, external, construct, conclusion, and reliability threats with mitigations
%% from docs/study_design.md §13.

\subsection{Internal validity}
\label{sec:threats:internal}

\subsection{External validity}
\label{sec:threats:external}

\subsection{Construct validity}
\label{sec:threats:construct}

\subsection{Conclusion validity}
\label{sec:threats:conclusion}

\subsection{Reliability and reproducibility}
\label{sec:threats:reliability}

\FloatBarrier
\section{Conclusion}
\label{sec:conclusion}

%% Goal: Summarize answers to RQ1--RQ7, restate contributions, future work (stretch tier,
%% repair pipeline, cross-model generalization). Numeric summaries: reported after campaign freeze.

\subsection{Summary of findings}
\label{sec:conclusion:summary}

\subsection{Contributions revisited}
\label{sec:conclusion:contributions}

\subsection{Future work}
\label{sec:conclusion:future}

%% Reported after campaign freeze.


\section*{Declarations}
\subsection*{Funding}
%% Goal: Funding sources (if any).
\subsection*{Conflict of interest}
%% Goal: COI statement.
\subsection*{Data availability}
%% Goal: Zenodo / repository DOI after M9.
\subsection*{Code availability}
%% Goal: behavioral-fsm-benchmark repository URL and license.

\bibliography{references/references}

\end{document}
